\scanpage{157}\mspage{12}%
nification des $\mathbb{E}\mathrm{xt}^{i}$ (hyperext), à un moment où Serre avait développé le formalisme des groupes proalgébriques, on n'avait pas trop peur de le prendre là-dedans -- au sens d'une catégorie dérivée qui à ce moment là n'avait encore jamais été définie et étudiée en forme.... (On a fait quand même un peu de chemin depuis !) J'ai bien l'impression, à la suite de ces antiques cogitations, pour heuristiques qu'elles furent, qu'on doit pouvoir développer à l'heure actuelle une telle théorie des $J_{*X/k}$, en coh fppf disons, donnant une formule $(*)$ sans limitation sur le dim $i$ de la cohomologie. (NB mais $J_*$ ne s'arrête évidemment plus en $\dim X=n$, mais en dim $2n$. Il est possible néanmoins que les composants $J_i$ soient de dim $0$ pour $i>n$).

Je crois bien que la construction de $J_*$ ne commute pas au chgt de base, mais que dans la cat. dérivée seulement elle y commute.

D) Soit $\deleted{f}X:/k$ schéma lisse \inserted{\deleted{\unsure{si|sic}}} (pour simplifier) sur corps $k$, séparé de \unsure{V.f.|T.f.|t.f.}, \inserted{de dim. relative $d$,} et $n$ entier $>0$. Si $n$ est premier à la caract. et si $F$ est un
% archive sheet number: 157
